\labsection{10}{File records and indexed block allocation}{sec:lab10}
\begin{labproject}{Separate application record formats from file-system block allocation}
\textbf{Suggested time:} 100--110 minutes \hfill
\textbf{Level:} intermediate\par
\textbf{Starter files:} \texttt{text\_records.c}, \texttt{binary\_records.c}, \texttt{indexed\_allocation.c}.\par
\textbf{Evidence to keep:} text/binary comparison, successful and missing search cases, and allocation-state trace.

\begin{labmode}
This practical deliberately works at two layers. First the application decides how records become bytes. Then a separate simulator shows how a file system could map a file through an index block to storage blocks.
\end{labmode}

\medskip\noindent\textbf{Learning outcomes.}
\begin{enumerate}
  \item compare text and raw binary record encodings;
  \item explain why sequential record search is simple but may require scanning many records;
  \item describe indexed allocation and explain why validate-then-commit protects allocation state.
\end{enumerate}

\medskip\noindent\textbf{Part A --- Text records are easy to inspect.}
\begin{lstlisting}[style=osbash]
cd experiments/Lab10
gcc -Wall -Wextra text_records.c -o text_records
./text_records
cat employees.txt
ls -l employees.txt
\end{lstlisting}
Identify the record delimiter and field delimiter. Change the target ID to a missing value, rerun, and ensure the program reports the missing record clearly.

\medskip\noindent\textbf{Part B --- Binary records preserve machine representation.}
\begin{lstlisting}[style=osbash]
gcc -Wall -Wextra binary_records.c -o binary_records
./binary_records
ls -l students.bin
hexdump -C students.bin | head
\end{lstlisting}
If \texttt{hexdump} is unavailable, use \texttt{od -Ax -tx1 students.bin | head}. Compare readability and size with the text file.

\begin{pitfall}
Writing a C structure directly with \texttt{fwrite()} is convenient for a controlled lab, but it is not a portable long-term interchange format. Structure padding, integer representation, endianness, and compiler choices can differ.
\end{pitfall}

\medskip\noindent\textbf{Part C --- Trace indexed allocation.}
\begin{lstlisting}[style=osbash]
gcc -Wall -Wextra indexed_allocation.c -o indexed_allocation
./indexed_allocation
\end{lstlisting}
Draw file A as one index block pointing to its three data blocks. Then explain why file B is rejected and verify that the printed allocation map is unchanged after that rejected request.

\medskip\noindent\textbf{Part D --- Explain validate first, commit second.}
Read the two functions \texttt{request\_is\_valid()} and \texttt{allocate\_file()}. In 4--6 sentences explain why checking the whole request before marking any block avoids a partially committed allocation after an error.

\medskip\noindent\textbf{Part E --- Connect the two layers.}
Draw two small diagrams: one showing fields inside an employee/student record, and one showing logical file blocks reached through an index block. Label which decisions belong to the application and which belong to the file system.

\begin{verification}
A complete submission includes a text/binary comparison, one successful and one missing record search, an indexed-allocation diagram, and evidence that a rejected overlapping allocation leaves block state unchanged.
\end{verification}

\begin{tryit}
Optional extension: add a deletion operation to the allocation simulator that frees an index block and all data blocks belonging to one simulated file. Print the allocation map before and after deletion.
\end{tryit}
\end{labproject}
